\subsection{DERIVATIVE OF A PRODUCT}

Let us now consider \(p\) topological vector spaces \(\mathrm{E}_i\) ( \(1 \leqslant i \leqslant p\) ) over \(\mathbf{R}\) , and a continuous multilinear \(^1\) map (which we shall denote by

\[
\left(\mathbf {x} _ {1}, \mathbf {x} _ {2}, \dots , \mathbf {x} _ {p}\right) \mapsto \left[ \mathbf {x} _ {1}. \mathbf {x} _ {2} \dots \mathbf {x} _ {p} \right])
\]

of \(E_{1} \times E_{2} \times \cdots \times E_{p}\) into a topological vector space F over R.

PROPOSITION 3. For each index \(i\) ( \(1 \leqslant i \leqslant p\) ) let \(\mathbf{f}_i\) be a function defined on an interval \(\mathrm{I} \subset \mathbf{R}\) , taking values in \(\mathrm{E}_i\) , and differentiable at the point \(x_0 \in \mathrm{I}\) . Then the function

\[
x \mapsto [ \mathbf {f} _ {1} (x). \mathbf {f} _ {2} (x) \dots \mathbf {f} _ {p} (x) ]
\]

defined on I with values in F has a derivative equal to

\[
\sum_ {i = 1} ^ {p} \left[ \mathbf {f} _ {1} (x _ {0}) \dots \mathbf {f} _ {i - 1} (x _ {0}). \mathbf {f} _ {i} ^ {\prime} (x _ {0}). \mathbf {f} _ {i + 1} (x _ {0}) \dots \mathbf {f} _ {p} (x _ {0}) \right]\tag{1}
\]

at \(x_0\) .

Let us put \(\mathbf{h}(x) = [\mathbf{f}_1(x).\mathbf{f}_2(x)\dots \mathbf{f}_p(x)]\) ; then, by the identity

\[
\left[ \mathbf {b} _ {1}. \mathbf {b} _ {2} \dots \mathbf {b} _ {p} \right] - \left[ \mathbf {a} _ {1}. \mathbf {a} _ {2} \dots \mathbf {a} _ {p} \right] = \sum_ {i = 1} ^ {p} \left[ \mathbf {b} _ {1} \dots \mathbf {b} _ {i - 1}. (\mathbf {b} _ {i} - \mathbf {a} _ {i}). \mathbf {a} _ {i + 1} \dots \mathbf {a} _ {p} \right],
\]

we can write

\[
\mathbf {h} (x) - \mathbf {h} (x _ {0}) = \sum_ {i = 1} ^ {p} \left[ \mathbf {f} _ {1} (x) \dots \mathbf {f} _ {i - 1} (x). (\mathbf {f} _ {i} (x) - \mathbf {f} _ {i} (x _ {0})). \mathbf {f} _ {i + 1} (x _ {0}) \dots \mathbf {f} _ {p} (x _ {0}) \right].
\]

On multiplying both sides by \(\frac{1}{x - x_0}\) and letting \(x\) approach \(x_0\) in I, we obtain the expression (1), since both the map

\[
(\mathbf {x} _ {1}, \mathbf {x} _ {2}, \dots , \mathbf {x} _ {p}) \mapsto [ \mathbf {x} _ {1}. \mathbf {x} _ {2} \dots \mathbf {x} _ {p} ]
\]

and addition in F are continuous.

\(^{1}\) Recall (Alg., II, p.~265) that a map f of \(E_{1} \times E_{2} \times \cdots \times E_{p}\) into F is said to be multilinear if each partial mapping

\[
\mathbf {x} _ {i} \mapsto \mathbf {f} (\mathbf {a} _ {1}, \dots , \mathbf {a} _ {i - 1}, \mathbf {x} _ {i}, \mathbf {a} _ {i + 1}, \dots , \mathbf {a} _ {p})
\]

from \(E_{i}\) into F ( \(1 \leqslant i \leqslant p\) ) is a linear map, the \(a_{j}\) for indices \(j \neq i\) being arbitrary in \(E_{j}\) . We note that if the \(E_{i}\) are finite dimensional over R then every multilinear map of \(E_{1} \times E_{2} \times \cdots \times E_{p}\) into F is necessarily continuous; this need not be so if some of these spaces are topological vector spaces of infinite dimension.

When some of the functions \(\mathbf{f}_i\) are constant, the terms in the expression (1) containing their derivatives \(\mathbf{f}_i'(x_0)\) are zero.

Let us consider in detail the particular case \(p = 2\) , the most important in applications: if \((\mathbf{x}, \mathbf{y}) \mapsto [\mathbf{x}.\mathbf{y}]\) is a continuous bilinear map of \(\mathrm{E} \times \mathrm{F}\) into \(\mathrm{G}\) , (E, F, G being topological vector spaces over \(\mathbf{R}\) ), and \(\mathbf{f}\) and \(\mathbf{g}\) are two vector functions, differentiable at \(x_0\) , with values in E and F respectively, then the vector function \(x \mapsto [\mathbf{f}(x).\mathbf{g}(x)]\) (which we denote by \([\mathbf{f}.\mathbf{g}]\) ) has a derivative equal to \([\mathbf{f}'(x_0).\mathbf{g}(x_0)] + [\mathbf{f}(x_0).\mathbf{g}'(x_0)]\) at \(x_0\) . In particular, if \(\mathbf{a}\) is a constant vector, then \([\mathbf{a}.\mathbf{f}]\) (resp. \([\mathbf{f}.\mathbf{a}]\) ) has a derivative equal to \([\mathbf{a}.\mathbf{f}'(x_0)]\) (resp. \([\mathbf{f}'(x_0).\mathbf{a}]\) ) at \(x_0\) .

If \(\mathbf{f}\) and \(\mathbf{g}\) are both differentiable on I then so is {[}f.g{]}, and we have

\[
\left[ \mathbf {f}. \mathbf {g} \right] ^ {\prime} = \left[ \mathbf {f} ^ {\prime}. \mathbf {g} \right] + \left[ \mathbf {f}. \mathbf {g} ^ {\prime} \right].\tag{2}
\]

Examples. 1) Let \(f\) be a real function, \(\mathbf{g}\) a vector function, both differentiable at a point \(x_0\) ; the function \(\mathbf{g}f\) has a derivative equal to \(\mathbf{g}'(x_0)f(x_0) + \mathbf{g}(x_0)f'(x_0)\) at \(x_0\) . In particular, if \(\mathbf{a}\) is constant, then \(\mathbf{a}f\) has derivative \(\mathbf{a}f'(x_0)\) . This last remark, in conjunction with example 1 of I, p.~5, proves that if \(\mathbf{f} = (f_i)_{1 \leqslant i \leqslant n}\) is a vector function with values in \(\mathbf{R}^n\) , then for \(\mathbf{f}\) to be differentiable at the point \(x_0\) it is necessary and sufficient that each of the real functions \(f_i\) ( \(1 \leqslant i \leqslant n\) ) be differentiable there: for, if \((\mathbf{e}_i)_{1 \leqslant i \leqslant n}\) is the canonical basis of \(\mathbf{R}^n\) , we can write \(\mathbf{f} = \sum_{i=1}^{n} \mathbf{e}_i f_i\) .

\begin{enumerate}
\def\labelenumi{\arabic{enumi})}
\setcounter{enumi}{1}
\tightlist
\item
  The real function \(x^{n}\) arises from the multilinear function
\end{enumerate}

\[
(x _ {1}, x _ {2}, \ldots , x _ {n}) \mapsto x _ {1} x _ {2} \ldots x _ {n}
\]

defined on \(R^{n}\) , by substituting x for each of the \(x_{i}\) ; so prop. 3 shows that \(x^{n}\) is differentiable on R and has derivative \(nx^{n-1}\) . As a result the polynomial function \(a_{0}x^{n} + a_{1}x^{n-1} + \cdots + a_{n-1}x + a_{n}\) (the \(a_{i}\) being constant vectors) has derivative

\[
n \mathbf {a} _ {0} x ^ {n - 1} + (n - 1) \mathbf {a} _ {1} x ^ {n - 2} + \dots + \mathbf {a} _ {n - 1};
\]

when the \(\mathbf{a}_i\) are real numbers this function coincides with the derivative of a polynomial function as defined in Algebra (A, IV).

\begin{enumerate}
\def\labelenumi{\arabic{enumi})}
\setcounter{enumi}{2}
\tightlist
\item
  The euclidean scalar product \((\mathbf{x}|\mathbf{y})\) (Gen.~Top., VI, p.~40) is a bilinear map (necessarily continuous) of \(R^{n}\times R^{n}\) into R. If f and g are two vector functions with values in \(R^{n}\) , and differentiable at the point \(x_{0}\) , then the real function \(x\mapsto(\mathbf{f}(x)|\mathbf{g}(x))\) has a derivative equal to \((\mathbf{f}'(x_{0})|\mathbf{g}(x_{0}))+(\mathbf{f}(x_{0})|\mathbf{g}'(x_{0}))\) at the point \(x_{0}\) . There is an analogous result for the hermitian scalar product on \(C^{n}\) , this space being considered as a vector space over R.
\end{enumerate}

Let us consider in particular the case where the euclidean norm \(\|\mathbf{f}(x)\|\) is constant, so that \((\mathbf{f}(x) \mid \mathbf{f}(x)) = \|\mathbf{f}(x)\|^{2}\) is also constant; on writing that the derivative of \((\mathbf{f}(x) \mid \mathbf{f}(x))\) vanishes at \(x_{0}\) we obtain \((\mathbf{f}(x_{0}) \mid \mathbf{f}'(x_{0})) = 0\) ; in other words, \(\mathbf{f}'(x_{0})\) is orthogonal to \(\mathbf{f}(x_{0})\) .

\begin{enumerate}
\def\labelenumi{\arabic{enumi})}
\setcounter{enumi}{3}
\item
  If E is a topological algebra over R (cf.~Introduction), the product xy of two elements of E is a continuous bilinear function of \((\mathbf{x}, \mathbf{y})\) ; if f and g have their values in E and are differentiable at the point \(x_{0}\) , then the function \(x \mapsto \mathbf{f}(x)\mathbf{g}(x)\) has a derivative equal to \(\mathbf{f}'(x_{0})\mathbf{g}(x_{0}) + \mathbf{f}(x_{0})\mathbf{g}'(x_{0})\) at \(x_{0}\) . In particular, if \(U(x) = (\alpha_{ij}(x))\) and \(V(x) = (\beta_{ij}(x))\) are two square matrices of order n, differentiable at \(x_{0}\) , their product UV has a derivative equal to \(U'(x_{0})V(x_{0}) + U(x_{0})V'(x_{0})\) at \(x_{0}\) (where \(U'(x) = (\alpha'_{ij}(x))\) and \(V'(x) = (\beta'_{ij}(x)))\) .
\item
  The determinant \(\operatorname{det}(\mathbf{x}_1, \mathbf{x}_2, \ldots, \mathbf{x}_n)\) of \(n\) vectors \(\mathbf{x}_i = (x_{ij})_{1 \leqslant j \leqslant n}\) from the space \(\mathbf{R}^n\) (Alg., III, p.~522) being a (continuous) multilinear function of the \(\mathbf{x}_i\) , one sees that if the \(n^2\) real functions \(f_{ij}\) are differentiable at \(x_0\) , then their determinant \(g(x) = \det(f_{ij}(x))\) has a derivative equal to
\end{enumerate}

\[
\sum_ {i = 1} ^ {n} \left[ \mathbf {f} _ {1} (x _ {0}), \dots , \mathbf {f} _ {i - 1} (x _ {0}), \mathbf {f} _ {i} ^ {\prime} (x _ {0}), \mathbf {f} _ {i + 1} (x _ {0}), \dots , \mathbf {f} _ {n} (x _ {0}) \right]
\]

at \(x_{0}\) , where \(\mathbf{f}_{i}(x)=(f_{ij}(x))_{1\leqslant j\leqslant n}\) ; in other words, one obtains the derivative of a determinant of order n by taking the sum of the n determinants formed by replacing, for each i, the terms of the \(i^{th}\) column by their derivatives.

Remark. If \(U(x)\) is a square matrix which is differentiable and invertible at the point \(x_{0}\) , then the derivative of its determinant \(\Delta(x) = \det(U(x))\) can be expressed through the derivative of \(U(x)\) by the formula

\[
\varDelta^ {\prime} (x _ {0}) = \varDelta (x _ {0}). \mathrm{Tr} (U ^ {\prime} (x _ {0}) U ^ {- 1} (x _ {0})).\tag{3}
\]

Indeed, let us put \(U(x_0 + h) = U(x_0) + hV\) ; then, by definition, \(V\) tends to \(U'(x_0)\) when \(h\) tends to 0. One can write

\[
\Delta (x _ {0} + h) = \Delta (x _ {0}). \det (I + h V U ^ {- 1} (x _ {0})).
\]

Now \(\det(I + hX) = 1 + h\operatorname{Tr}(X) + \sum_{k=2}^{n} \lambda_k h^k\) , the \(\lambda_k (k \geqslant 2)\) being polynomials in the elements of the matrix X; since the elements of \(VU^{-1}(x_0)\) have a limit when h tends to 0, we indeed obtain the formula (3).

